\section{Baselines}
\label{sec:baselines}

A baseline is a model that requires no fitting: no parameters are
estimated from data. Its purpose is to define the bar every learned model
must clear (Section~\ref{sec:metrics}'s skill score is measured against
it).

\subsection{Persistence}

\begin{equation}
\label{eq:persistence}
\hat y_{c,t+1} = y_{c,t}
\end{equation}

Equation~\ref{eq:persistence} predicts that next year's level equals this
year's level. This is the naive forecast, and it is the optimal forecast
under a random-walk-without-drift model of the underlying series
\cite{hyndman2021forecasting}. It is a strong baseline for country-level
CO$_2$ specifically because the underlying physical stock (power plants,
vehicle fleets, industrial capital) changes slowly: the median absolute
year-over-year change in this dataset since 2010 is approximately 4.4\%
of the level (Session 1 EDA).

\subsection{Linear Trend (Drift)}

\begin{equation}
\label{eq:trend}
\hat y_{c,t+1} = y_{c,t} + \bar\Delta_{c,t}^{(5)},
\qquad
\bar\Delta_{c,t}^{(5)} = \frac{y_{c,t} - y_{c,t-5}}{5}
\end{equation}

Equation~\ref{eq:trend} extrapolates the average yearly change over the
preceding five years one step forward. It is this project's
five-year-windowed variant of the classical drift method
\cite{hyndman2021forecasting}, which in its general form draws a line
between a series' first and last observed points and extrapolates it. The
trade-off is explicit: Equation~\ref{eq:trend} outperforms
Equation~\ref{eq:persistence} for a country on a sustained trajectory, and
underperforms it the moment a country's direction reverses, since the
linear extrapolation has no mechanism to anticipate a turn.

\subsection{Bias direction}

For any series with a genuine sustained trend, persistence
(Equation~\ref{eq:persistence}) is a biased estimator of $y_{c,t+1}$: if
$y_{c,t}$ is rising, persistence systematically under-predicts; if
$y_{c,t}$ is falling (for example, a country whose emissions peaked and
have since declined), persistence systematically over-predicts. The sign
of the bias tracks the sign of the country's own recent trend, not a
fixed direction.
